% from ega4-2-fr.tex
\oldpage[IV-2]{189}

parce que $A'$ est universellement japonais (7.6.5)). Pour cette raison,
les conditions $b'_1)$ et c) sont aussi équivalentes en vertu de (5.11.4).
\end{proof}

\begin{proposition}[7.2.4]
\label{IV.2.s7.2.4-fr}
Soit $A$ un anneau local noethérien et posons
$X=\operatorname{Spec}(A)$. Les conditions suivantes sont équivalentes:
\begin{enumerate}
\item[a)] Pour tout anneau quotient intègre $B$ de $A$, l'anneau $B^{(1)}$
est une $B$-algèbre finie.
\item[b)] Pour tout $\mathscr O_X$-Module cohérent $\mathscr F$ et toute
partie $Z$, stable par spécialisation, et telle que pour tout
$x\in\operatorname{Ass}(\mathscr F)\cap(X-Z)$, on ait
$\operatorname{codim}(\overline{\{x\}}\cap Z,\overline{\{x\}})\geqslant2$,
le $\mathscr O_X$-Module $\mathscr H^0_{X/Z}(\mathscr F)$ est cohérent.
\item[c)] Pour toute partie fermée $T$ de $X$ et tout
$\mathscr O_U$-Module cohérent $\mathscr G$ (où $U=X-T$) tel que, pour tout
$x\in\operatorname{Ass}(\mathscr G)$, on ait
$\operatorname{codim}(\overline{\{x\}}\cap T,\overline{\{x\}})\geqslant2$,
$i_*(\mathscr G)$ (où $i:U\to X$ est l'injection canonique) est un
$\mathscr O_X$-Module cohérent.
\item[d)] Pour tout anneau quotient intègre $B$ de $A$, et tout idéal
$\mathfrak J$ de hauteur $\geqslant2$ dans $B$, l'anneau
$\displaystyle\bigcap_{\mathfrak p\not\supset\mathfrak J}B_{\mathfrak p}$
est une $B$-algèbre finie.
\item[e)] Pour tout anneau quotient intègre $B$ de $A$ et tout anneau local
$C=B_{\mathfrak q}$ en un idéal premier $\mathfrak q$ de $B$, tel que
$\dim(C)\geqslant2$, l'anneau $C^{(\omega)}$ est une $C$-algèbre finie (ou,
ce qui revient au même, si $Y=\operatorname{Spec}(C)$ et si $U$ est le
complémentaire dans $Y$ du point fermé de $C$ et $j:U\to Y$ l'injection
canonique, $j_*(\mathscr O_U)$ est un $\mathscr O_Y$-Module cohérent).
\item[f)] Pour tout anneau quotient intègre $B$ de $A$, tout anneau local
$C=B_{\mathfrak q}$ en un idéal premier $\mathfrak q$ de $B$, tel que
$\dim(C)\geqslant2$, et pour tout idéal
$\mathfrak r\in\operatorname{Ass}(\widehat C)$, on a
$\dim(\widehat C/\mathfrak r)\geqslant2$.
\end{enumerate}
\end{proposition}

\begin{proof}
On sait déjà (5.11.6) que a) et b) sont équivalentes, ainsi que c) et d),
et que a) entraîne d). L'équivalence de a) et d) dans le cas actuel résulte
de (7.2.3) appliqué à un anneau quotient intègre $B$ de $A$, la condition
d) étant une formulation équivalente de la condition (7.2.3, b)) en vertu
de (5.9.3.1). L'équivalence de e) et f) est une conséquence de (7.2.2)
appliqué au $\mathscr O_Y$-Module cohérent $\mathscr O_Y$ lui-même. On a
déjà remarqué (5.11.7, (ii)) que si $A$ vérifie a), il en est de même de
toute $A$-algèbre finie et de tout anneau de fractions de $A$. La condition
a) pour $A$ entraîne donc (avec les notations de e)) que l'anneau $C$
vérifie a), donc aussi c); mais comme $C$ est intègre et $\dim(C)\geqslant
2$, on peut appliquer à $C$ la condition c) en prenant pour $T$ le point
fermé de $Y=\operatorname{Spec}(C)$; cela prouve que a) entraîne e). Reste
à prouver que e) entraîne a); on peut évidemment se borner au cas où $A$
est intègre et $B=A$, et il s'agit alors de montrer que la condition e)
entraîne la condition (7.2.3, c)). Avec les notations de (7.2.3, c)),
considérons donc un point
$y'\in f^{-1}(T)\cap\overline{\{x'\}}$, et posons $y=f(y')$ et
$C=\mathscr O_{X,y}$; comme $y\in T$, on a par hypothèse
$\dim(C)\geqslant2$, et par suite (avec les notations de e)),
$j_*(\mathscr O_U)$ est un $\mathscr O_Y$-Module cohérent. Or, posons
$Y'=X'\times_XY$; le morphisme $f:X'\to X$ étant plat, il en est de même de
$g=f_{(Y)}:Y'\to Y$. D'ailleurs l'espace sous-jacent à $Y'$ s'identifie à
$f^{-1}(Y)$ (I, 3.6.5), et comme $A$ est intègre et $f$ plat,
$\operatorname{Ass}(\mathscr O_{X'})$ est contenu dans la fibre du point
générique de $X$ (3.3.2); ce dernier étant contenu dans $Y$, on a
$x'\in Y'$. Soit $U'=g^{-1}(U)$, et soit $j'$ l'injection canonique
$U'\to Y'$; comme $j_*(\mathscr O_U)$ est un $\mathscr O_Y$-Module
cohérent et que $g$ est plat, il résulte de (5.9.4) que
$j'_*(\mathscr O_{U'})$ est un $\mathscr O_{Y'}$-Module cohérent. Or on a
$x'\in\operatorname{Ass}(\mathscr O_{U'})$, donc on conclut de (5.10.10)
que l'on a, dans $Y'$,

